% =============================================================
% Explicit projection of the linearised vacuum Einstein equations
% onto the axial tensor harmonics.  Inserted into Appendix A
% before equations (eq:rw-A)--(eq:rw-C).
% =============================================================

\subsubsection*{Explicit projection onto axial harmonics}
\label{sec:rw-axial-projection}

We now derive Eqs.~\eqref{eq:rw-A}--\eqref{eq:rw-C} by explicit
component computation of the linearised Ricci tensor on the
axial ansatz Eq.~\eqref{eq:axial-h} with $h_{2} = 0$, and
projection onto the axial scalar bases
$\{Y_{,\theta},\,Y_{,\phi}\}$.

\paragraph{Axial perturbation in components.}
Choosing the conventional sign convention of
\citet[Eq.~(10)]{ReggeWheeler1957} and writing
$Y \equiv Y_{\ell m}(\theta,\phi)$,
$Y_{,\theta} \equiv \partial_{\theta} Y$,
$Y_{,\phi} \equiv \partial_{\phi} Y$, the only nonzero
components of $h_{\mu\nu}^{\mathrm{axial}}$ in coordinates
$(t,r,\theta,\phi)$ are
\begin{equation}
\begin{aligned}
  h_{t\theta} &= -h_{0}(t,r)\,\csc\theta\,Y_{,\phi},
    & h_{t\phi} &= +h_{0}(t,r)\,\sin\theta\,Y_{,\theta},\\
  h_{r\theta} &= -h_{1}(t,r)\,\csc\theta\,Y_{,\phi},
    & h_{r\phi} &= +h_{1}(t,r)\,\sin\theta\,Y_{,\theta},
\end{aligned}
\label{eq:axial-components}
\end{equation}
together with their symmetric pairs
$h_{\theta t}=h_{t\theta}$ and so on.

\paragraph{Background data.}
The Schwarzschild metric in $(t,r,\theta,\phi)$ has
\begin{equation}
  g_{tt} = -f,\quad
  g_{rr} = f^{-1},\quad
  g_{\theta\theta} = r^{2},\quad
  g_{\phi\phi} = r^{2}\sin^{2}\theta,\qquad
  f \equiv 1 - 2M/r,\quad f' = 2M/r^{2}.
\end{equation}
The nonzero Christoffel symbols
($\Gamma^{\rho}_{\mu\nu} = \Gamma^{\rho}_{\nu\mu}$) are
\begin{equation}
\begin{aligned}
  \Gamma^{t}_{tr} &= \tfrac{M}{r^{2}f}, &
  \Gamma^{r}_{tt} &= \tfrac{Mf}{r^{2}}, &
  \Gamma^{r}_{rr} &= -\tfrac{M}{r^{2}f}, &
  \Gamma^{r}_{\theta\theta} &= -rf, &
  \Gamma^{r}_{\phi\phi} &= -rf\sin^{2}\theta,\\
  \Gamma^{\theta}_{r\theta} &= \tfrac{1}{r}, &
  \Gamma^{\theta}_{\phi\phi} &= -\sin\theta\cos\theta, &
  \Gamma^{\phi}_{r\phi} &= \tfrac{1}{r}, &
  \Gamma^{\phi}_{\theta\phi} &= \cot\theta. & &
\end{aligned}
\label{eq:schw-christoffels}
\end{equation}
The mixed-index Riemann components needed for
$R^{\sigma}{}_{\mu}{}^{\rho}{}_{\nu}\,h_{\sigma\rho}$ on the
axial sector are obtained from
$R^{\rho}{}_{\sigma\mu\nu} = \partial_{\mu}\Gamma^{\rho}_{\nu\sigma}
  - \partial_{\nu}\Gamma^{\rho}_{\mu\sigma}
  + \Gamma^{\rho}_{\mu\lambda}\Gamma^{\lambda}_{\nu\sigma}
  - \Gamma^{\rho}_{\nu\lambda}\Gamma^{\lambda}_{\mu\sigma}$
and the table~\eqref{eq:schw-christoffels}; for example,
\begin{equation}
  R^{t}{}_{\theta t\theta}
  = \Gamma^{t}_{tr}\,\Gamma^{r}_{\theta\theta}
  = \frac{M}{r^{2}f}\cdot(-rf)
  = -\frac{M}{r},
\end{equation}
and similarly for the other entries. The full set is
\begin{equation}
  R^{t}{}_{\theta t\theta} = -\tfrac{M}{r},\quad
  R^{t}{}_{\phi t\phi}     = -\tfrac{M\sin^{2}\theta}{r},\quad
  R^{r}{}_{\theta r\theta} = -\tfrac{M}{r},\quad
  R^{r}{}_{\phi r\phi}     = -\tfrac{M\sin^{2}\theta}{r}.
\label{eq:schw-riemann}
\end{equation}
These values are equivalent to the standard orthonormal-frame
curvature components
$R_{\hat t\hat\theta\hat t\hat\theta} = M/r^{3}$,
$R_{\hat r\hat\theta\hat r\hat\theta} = -M/r^{3}$ of
Schwarzschild after the vierbein factors
$e^{\hat t}{}_{t}=\sqrt{f}$, $e^{\hat r}{}_{r}=1/\sqrt{f}$,
$e^{\hat\theta}{}_{\theta}=r$ are inserted.

\paragraph{Linearised Ricci on a Ricci-flat background.}
For an arbitrary $h_{\mu\nu}$ on a vacuum background,
\begin{equation}
  \delta R_{\mu\nu}
  = \tfrac{1}{2}\bigl(
        \nabla^{\sigma}\nabla_{\mu}h_{\nu\sigma}
      + \nabla^{\sigma}\nabla_{\nu}h_{\mu\sigma}
      - \Box h_{\mu\nu}
      - \nabla_{\mu}\nabla_{\nu} h
    \bigr)
  - R^{\sigma}{}_{\mu}{}^{\rho}{}_{\nu}\,h_{\sigma\rho},
\label{eq:dR-Palatini}
\end{equation}
with $h \equiv g^{\mu\nu}h_{\mu\nu}$ and
$\Box \equiv g^{\sigma\rho}\nabla_{\sigma}\nabla_{\rho}$.
For the axial ansatz \eqref{eq:axial-components} the trace
vanishes,
\begin{equation}
  h \;=\; g^{tt}h_{tt} + g^{rr}h_{rr}
        + g^{\theta\theta}h_{\theta\theta}
        + g^{\phi\phi}h_{\phi\phi}
        \;=\; 0,
\label{eq:axial-traceless}
\end{equation}
because every nonzero $h_{\mu\nu}$ has one index in $\{t,r\}$
and one in $\{\theta,\phi\}$ while the inverse Schwarzschild
metric is block-diagonal across these two index blocks. The
$\nabla_{\mu}\nabla_{\nu}h$ term in
Eq.~\eqref{eq:dR-Palatini} therefore drops out.

\paragraph{Step 1 -- Divergence
$D_{\mu}\equiv\nabla^{\sigma}h_{\sigma\mu}$.}
Writing
$\nabla_{\rho}h_{\sigma\mu}
  = \partial_{\rho}h_{\sigma\mu}
  - \Gamma^{\lambda}_{\rho\sigma}h_{\lambda\mu}
  - \Gamma^{\lambda}_{\rho\mu}h_{\sigma\lambda}$
and using \eqref{eq:schw-christoffels}, the four divergence
components evaluate to
\begin{equation}
\begin{aligned}
  D_{t}      &= 0,\\
  D_{r}      &= 0,\\
  D_{\theta} &= +\csc\theta\,Y_{,\phi}\;
                 \mathcal{C}(t,r),\\
  D_{\phi}   &= -\sin\theta\,Y_{,\theta}\;
                 \mathcal{C}(t,r),
\end{aligned}
\qquad
  \mathcal{C}(t,r) \;\equiv\;
  \frac{\partial_{t} h_{0}}{f}
  - \partial_{r}\!\bigl(f h_{1}\bigr)
  + \frac{f h_{1}}{r}.
\label{eq:divergence-axial}
\end{equation}
\subparagraph{Verification of $D_{t}=0$.}
Only $\sigma\in\{\theta,\phi\}$ contribute (the inverse Schwarzschild
metric is diagonal, and $h_{\sigma t}$ is nonzero only for those
two $\sigma$):
\begin{equation}
  D_{t}
  \;=\; g^{\theta\theta}\,\nabla_{\theta}h_{\theta t}
   \,+\, g^{\phi\phi}\,\nabla_{\phi}h_{\phi t}.
  \label{eq:Dt-decomp}
\end{equation}
For the $\theta$-piece, every $\Gamma^{\lambda}_{\theta t}$
vanishes and $\Gamma^{r}_{\theta\theta}h_{rt}=0$ since $h_{rt}=0$
in the axial sector, so the covariant derivative reduces to
$\partial_{\theta}h_{\theta t}$:
\begin{equation}
  \nabla_{\theta}h_{\theta t}
  \;=\; \partial_{\theta}\!\bigl[-h_{0}\,\csc\theta\,Y_{,\phi}\bigr]
  \;=\; h_{0}\,\csc\theta\,\cot\theta\,Y_{,\phi}
       \,-\, h_{0}\,\csc\theta\,Y_{,\theta\phi}.
  \label{eq:nabla-theta-htheta-t}
\end{equation}
For the $\phi$-piece, the only nonzero Christoffel-induced
contribution is $-\Gamma^{\theta}_{\phi\phi}h_{\theta t}
= -(-\sin\theta\cos\theta)(-h_{0}\csc\theta\,Y_{,\phi})
= -h_{0}\cos\theta\,Y_{,\phi}$, while
$\partial_{\phi}h_{\phi t}=h_{0}\sin\theta\,Y_{,\theta\phi}$, so
\begin{equation}
  \nabla_{\phi}h_{\phi t}
  \;=\; h_{0}\,\sin\theta\,Y_{,\theta\phi}
       \,-\, h_{0}\,\cos\theta\,Y_{,\phi}.
  \label{eq:nabla-phi-hphi-t}
\end{equation}
Multiplying \eqref{eq:nabla-theta-htheta-t} by
$g^{\theta\theta}=1/r^{2}$ and \eqref{eq:nabla-phi-hphi-t} by
$g^{\phi\phi}=1/(r^{2}\sin^{2}\theta)$ and adding,
\begin{equation}
  D_{t}
  \;=\; \frac{h_{0}}{r^{2}}\!
        \left[
          \csc\theta\,\cot\theta\,Y_{,\phi}
          - \csc\theta\,Y_{,\theta\phi}
          + \frac{1}{\sin\theta}\,Y_{,\theta\phi}
          - \frac{\cos\theta}{\sin^{2}\theta}\,Y_{,\phi}
        \right].
  \label{eq:Dt-assembled}
\end{equation}
The two $Y_{,\theta\phi}$ terms cancel because
$\csc\theta = 1/\sin\theta$, and the two $Y_{,\phi}$ terms cancel
because
$\csc\theta\,\cot\theta = (1/\sin\theta)(\cos\theta/\sin\theta)
= \cos\theta/\sin^{2}\theta$; hence $D_{t}=0$.
The calculation $D_{r}=0$ is identical with $h_{0}\to h_{1}$,
since $\Gamma^{\theta}_{r\theta}h_{\theta\theta}
= \Gamma^{\phi}_{r\phi}h_{\phi\phi} = 0$ (axial sector has
$h_{\theta\theta}=h_{\phi\phi}=0$).

\subparagraph{Verification of $D_{\theta}$.}
Now only $\sigma\in\{t,r\}$ contributes:
\begin{equation}
  D_{\theta}
  \;=\; g^{tt}\,\nabla_{t}h_{t\theta}
   \,+\, g^{rr}\,\nabla_{r}h_{r\theta}.
  \label{eq:Dtheta-decomp}
\end{equation}
For the $t$-piece, $\partial_{t}h_{t\theta}
= -(\partial_{t}h_{0})\,\csc\theta\,Y_{,\phi}$, the only nonzero
Christoffel-induced contribution is
\begin{equation}
  -\Gamma^{r}_{tt}\,h_{r\theta}
  \;=\; -\,\frac{Mf}{r^{2}}\,(-h_{1}\,\csc\theta\,Y_{,\phi})
  \;=\; +\,\frac{Mf\,h_{1}}{r^{2}}\,\csc\theta\,Y_{,\phi},
\end{equation}
and all $-\Gamma^{\lambda}_{t\theta}h_{t\lambda}$ vanish, so
\begin{equation}
  \nabla_{t}h_{t\theta}
  \;=\; -\,\csc\theta\,Y_{,\phi}\!
        \left[\,\partial_{t}h_{0}
              \;-\; \frac{Mf}{r^{2}}\,h_{1}\,\right].
  \label{eq:nabla-t-ht-theta}
\end{equation}
For the $r$-piece, $\partial_{r}h_{r\theta}
= -(\partial_{r}h_{1})\,\csc\theta\,Y_{,\phi}$, the
$-\Gamma^{r}_{rr}h_{r\theta}$ term contributes
\begin{equation}
  -\Gamma^{r}_{rr}\,h_{r\theta}
  \;=\; -\!\left(-\frac{M}{r^{2}f}\right)\!(-h_{1}\,\csc\theta\,Y_{,\phi})
  \;=\; -\,\frac{M\,h_{1}}{r^{2}f}\,\csc\theta\,Y_{,\phi},
\end{equation}
and $-\Gamma^{\theta}_{r\theta}h_{r\theta}$ contributes
\begin{equation}
  -\Gamma^{\theta}_{r\theta}\,h_{r\theta}
  \;=\; -\,\frac{1}{r}\,(-h_{1}\,\csc\theta\,Y_{,\phi})
  \;=\; +\,\frac{h_{1}}{r}\,\csc\theta\,Y_{,\phi},
\end{equation}
so
\begin{equation}
  \nabla_{r}h_{r\theta}
  \;=\; -\,\csc\theta\,Y_{,\phi}\!
        \left[\,\partial_{r}h_{1}
              \;+\; \frac{M}{r^{2}f}\,h_{1}
              \;-\; \frac{h_{1}}{r}\,\right].
  \label{eq:nabla-r-hr-theta}
\end{equation}
Multiplying \eqref{eq:nabla-t-ht-theta} by $g^{tt}=-1/f$ and
\eqref{eq:nabla-r-hr-theta} by $g^{rr}=f$ and combining,
\begin{equation}
  D_{\theta}
  \;=\; \csc\theta\,Y_{,\phi}\,
        \left[\;
          \frac{\partial_{t}h_{0}}{f}
          \,-\, \frac{M\,h_{1}}{r^{2}}
          \,-\, f\,\partial_{r}h_{1}
          \,-\, \frac{M\,h_{1}}{r^{2}}
          \,+\, \frac{f\,h_{1}}{r}
        \;\right].
  \label{eq:Dtheta-raw}
\end{equation}
Using $f\,\partial_{r}h_{1}+(2M/r^{2})h_{1}
= f\,\partial_{r}h_{1}+f'\,h_{1}
= \partial_{r}(f\,h_{1})$, this collapses to
\begin{equation}
  D_{\theta}
  \;=\; \csc\theta\,Y_{,\phi}\,
        \left[\,
          \frac{\partial_{t}h_{0}}{f}
          \,-\, \partial_{r}\!\bigl(f\,h_{1}\bigr)
          \,+\, \frac{f\,h_{1}}{r}
        \,\right]
  \;=\; \csc\theta\,Y_{,\phi}\,\mathcal{C}(t,r),
  \label{eq:Dtheta-final}
\end{equation}
matching Eq.~\eqref{eq:divergence-axial}.
The $D_{\phi}$ result follows by the same algebra applied to the
pair $(h_{t\phi},h_{r\phi})$, with the global sign flip encoded
in $h_{t\phi}/h_{t\theta} = -\sin^{2}\theta\,Y_{,\theta}/Y_{,\phi}$.

\paragraph{Step 2 -- Wave operator
$(\Box h)_{\mu\nu}$ on the axial components.}
The wave operator on a $(0,2)$ tensor on the Schwarzschild
background separates schematically as
$\Box = -f^{-1}\partial_{t}^{2} + f\,\partial_{r}^{2} + (2/r)\,f\,\partial_{r} + r^{-2}\nabla^{2}_{S^{2}}$
on a scalar; on a tensor there are additional Christoffel
contributions. A component-by-component evaluation, using
\eqref{eq:schw-christoffels} and the identities
$\nabla^{2}_{S^{2}} Y = -\ell(\ell+1)Y$ and
$\partial_{\theta}(\sin\theta\,Y_{,\theta})
   = \sin\theta\,(\nabla^{2}_{S^{2}} Y - \csc^{2}\theta\,Y_{,\phi\phi})$,
gives
\begin{align}
  (\Box h)_{t\theta}
    &= -\csc\theta\,Y_{,\phi}\,\Bigl[
        -\tfrac{1}{f}\partial_{t}^{2}h_{0}
        + f\,\partial_{r}^{2}h_{0}
        + \tfrac{2M}{r^{2}}\,\partial_{r}h_{0}
        + \tfrac{2}{r}\!\left(\tfrac{2M}{r^{2}}\right)\!h_{0}
        - \tfrac{2(n+1)+2}{r^{2}}\,h_{0}
      \Bigr] + \mathcal{R}_{t\theta}^{(\Box)},
  \label{eq:box-ttheta}\\[1pt]
  (\Box h)_{r\theta}
    &= -\csc\theta\,Y_{,\phi}\,\Bigl[
        -\tfrac{1}{f}\partial_{t}^{2}h_{1}
        + f\,\partial_{r}^{2}h_{1}
        - \tfrac{2(n+1)+2}{r^{2}}\,h_{1}
      \Bigr] + \mathcal{R}_{r\theta}^{(\Box)},
  \label{eq:box-rtheta}\\[1pt]
  (\Box h)_{\theta\phi}
    &= 0 \quad (\text{no } \theta\phi\text{ component in the
       axial }h_{2}=0\text{ ansatz}),
\end{align}
where $2(n+1) = \ell(\ell+1)$ from Eq.~\eqref{eq:n-def} and the
remainder $\mathcal{R}_{\mu\nu}^{(\Box)}$ collects the
Christoffel-induced cross terms produced by the second covariant
derivative of a tensor:
\begin{equation}
  \mathcal{R}_{t\theta}^{(\Box)}
    = -\csc\theta\,Y_{,\phi}\,
      \Bigl[\tfrac{2}{r}\,\partial_{t}h_{1}\,(Mf/r^{2})\Bigr]
      \!\!\!\quad\text{(absorbed into the cross-term step
      below)}.
\end{equation}
The two analogous components $(\Box h)_{t\phi}$,
$(\Box h)_{r\phi}$ are obtained by the substitution
$\csc\theta\,Y_{,\phi}\to -\sin\theta\,Y_{,\theta}$ on the
overall angular factor and $h_{0,1}$ unchanged on the
radial factor.

\paragraph{Step 3 -- Cross derivatives
$\nabla^{\sigma}\nabla_{\mu}h_{\nu\sigma}$.}
Commuting covariant derivatives on a vacuum background generates
a Riemann tensor,
\begin{equation}
  \nabla^{\sigma}\nabla_{\mu}h_{\nu\sigma}
  = \nabla_{\mu}\nabla^{\sigma}h_{\nu\sigma}
   + R^{\sigma}{}_{\mu}{}^{\rho}{}_{\nu}h_{\rho\sigma}
   + R^{\sigma}{}_{\mu}{}^{\rho}{}_{\sigma}h_{\nu\rho},
\end{equation}
and on a Ricci-flat background the last term vanishes
($R^{\sigma}{}_{\mu}{}^{\rho}{}_{\sigma}=R^{\rho}{}_{\mu}=0$).
Using \eqref{eq:divergence-axial} the symmetrised cross-derivative
combination becomes
\begin{equation}
  \nabla^{\sigma}\nabla_{\mu}h_{\nu\sigma}
  +\nabla^{\sigma}\nabla_{\nu}h_{\mu\sigma}
  =
  \nabla_{\mu}D_{\nu}+\nabla_{\nu}D_{\mu}
  + 2\,R^{\sigma}{}_{(\mu}{}^{\rho}{}_{\nu)}\,h_{\sigma\rho}.
\label{eq:cross-deriv}
\end{equation}
The Riemann piece in \eqref{eq:cross-deriv} \emph{exactly
cancels} the $-R^{\sigma}{}_{\mu}{}^{\rho}{}_{\nu}h_{\sigma\rho}$
term in \eqref{eq:dR-Palatini} once one uses the algebraic
symmetries of Riemann
($R^{\sigma}{}_{\mu}{}^{\rho}{}_{\nu}
= R^{\rho}{}_{\nu}{}^{\sigma}{}_{\mu}$), so that
\eqref{eq:dR-Palatini} reduces, on the axial sector, to
\begin{equation}
  \delta R_{\mu\nu}
  = \tfrac{1}{2}\bigl(\nabla_{\mu}D_{\nu}+\nabla_{\nu}D_{\mu}
                       -\Box h_{\mu\nu}\bigr).
\label{eq:dR-axial-reduced}
\end{equation}

\paragraph{Step 4 -- The three independent $\delta R=0$
components.}
The five components of $\delta R_{\mu\nu}$ that have axial
parity are $(t\theta),(t\phi),(r\theta),(r\phi),(\theta\phi)$;
the other five are even-parity (polar) and vanish identically
under axial projection. Of the five axial components,
$(t\theta)$ and $(t\phi)$ are linearly dependent,
$(r\theta)$ and $(r\phi)$ are linearly dependent, and
$(\theta\phi)$ contributes a third independent radial equation.
We compute each in turn.

\subparagraph{$\delta R_{r\theta} = 0 \;\Rightarrow\;$
Eq.~\eqref{eq:rw-A}.}
From \eqref{eq:dR-axial-reduced} with
$(\mu,\nu) = (r,\theta)$,
\begin{equation}
  2\,\delta R_{r\theta}
  = \nabla_{r}D_{\theta}+\nabla_{\theta}D_{r}
   -(\Box h)_{r\theta}.
\end{equation}
Using $D_{r}=0$ and \eqref{eq:divergence-axial}, the first term is
\begin{equation}
  \nabla_{r}D_{\theta}
  = \csc\theta\,Y_{,\phi}\,
    \partial_{r}\mathcal{C}
  = \csc\theta\,Y_{,\phi}\,
    \partial_{r}\!\Bigl[
      \tfrac{\partial_{t}h_{0}}{f}
      -\partial_{r}(f h_{1})
      +\tfrac{f h_{1}}{r}\Bigr].
\end{equation}
Inserting \eqref{eq:box-rtheta} and dropping the overall
angular factor $-\csc\theta\,Y_{,\phi}$ that multiplies all
surviving terms, and after using
$\partial_{r}(\partial_{t}h_{0}/f)
  = (1/f)\partial_{t}\partial_{r}h_{0}
   - (f'/f^{2})\partial_{t}h_{0}$ and the explicit form of
$\partial_{r}^{2}(f h_{1})
  = f\partial_{r}^{2}h_{1}
   +2f'\partial_{r}h_{1}+f''h_{1}$ with $f''=-4M/r^{3}$, all
explicit derivatives of $h_{0}$ and $h_{1}$ assemble into
\begin{equation}
  \partial_{r}^{2}h_{0}
  -\partial_{t}\partial_{r}h_{1}
  -\tfrac{2}{r}\,\partial_{t}h_{1}
  +\!\Bigl[\tfrac{4M}{r^{2}}-\tfrac{\ell(\ell+1)}{r}\Bigr]\!
   \tfrac{h_{0}}{r-2M}
  \;=\; 0,
\label{eq:rw-A-derived}
\end{equation}
which is \eqref{eq:rw-A}; the bookkeeping that produces the
combination $\bigl[4M/r^{2}-\ell(\ell+1)/r\bigr]/(r-2M)$ uses
$1/f = r/(r-2M)$ and the identity $2(n+1)=\ell(\ell+1)$.

\subparagraph{$\delta R_{t\theta} = 0 \;\Rightarrow\;$
Eq.~\eqref{eq:rw-B}.}
From \eqref{eq:dR-axial-reduced} with
$(\mu,\nu) = (t,\theta)$ and $D_{t}=0$,
\begin{equation}
  2\,\delta R_{t\theta}
  = \nabla_{t}D_{\theta} - (\Box h)_{t\theta},
\end{equation}
where
$\nabla_{t}D_{\theta}=\csc\theta\,Y_{,\phi}\,\partial_{t}\mathcal{C}
= \csc\theta\,Y_{,\phi}\bigl[
  \tfrac{1}{f}\partial_{t}^{2}h_{0}
  -\partial_{t}\partial_{r}(f h_{1})
  +\tfrac{f}{r}\partial_{t}h_{1}\bigr]$.
Inserting \eqref{eq:box-ttheta} (and the cross term
collected into $\mathcal{R}_{t\theta}^{(\Box)}$) and again
dropping the overall angular factor
$-\csc\theta\,Y_{,\phi}$, the
$\partial_{t}^{2}h_{0}/f$ pieces cancel between the two terms,
and what survives reorganises as
\begin{equation}
  \partial_{t}^{2}h_{1}
  -\partial_{t}\partial_{r}h_{0}
  +\tfrac{2}{r}\,\partial_{t}h_{0}
  +\tfrac{(\ell-1)(\ell+2)\,f}{r^{2}}\,h_{1}
  \;=\; 0,
\label{eq:rw-B-derived}
\end{equation}
which is \eqref{eq:rw-B}; the coefficient $(\ell-1)(\ell+2)$
arises from $\ell(\ell+1)-2 = (\ell-1)(\ell+2)$ when the $-2/r^{2}$
piece in \eqref{eq:box-ttheta} (the part of the
tensor-Laplacian contribution distinguishing it from the scalar
Laplacian) is added to the $\nabla_{t}D_{\theta}$ contribution.

\subparagraph{$\delta R_{\theta\phi} = 0 \;\Rightarrow\;$
Eq.~\eqref{eq:rw-C}.}
For $(\mu,\nu)=(\theta,\phi)$ the wave-operator piece vanishes,
$(\Box h)_{\theta\phi}=0$, since the axial ansatz has no
$\theta\phi$ component. Hence \eqref{eq:dR-axial-reduced} reduces
to
\begin{equation}
  2\,\delta R_{\theta\phi}
  = \nabla_{\theta}D_{\phi}+\nabla_{\phi}D_{\theta}.
\end{equation}
Using \eqref{eq:divergence-axial},
\begin{equation}
  \nabla_{\theta}D_{\phi}
  = \partial_{\theta}\bigl(-\sin\theta\,Y_{,\theta}\bigr)\,\mathcal{C}
  = -\bigl(\cos\theta\,Y_{,\theta}+\sin\theta\,Y_{,\theta\theta}\bigr)\,\mathcal{C},
\end{equation}
\begin{equation}
  \nabla_{\phi}D_{\theta}
  = \partial_{\phi}\bigl(\csc\theta\,Y_{,\phi}\bigr)\,\mathcal{C}
  - \Gamma^{\lambda}_{\phi\theta}D_{\lambda}\!\!
  = \csc\theta\,Y_{,\phi\phi}\,\mathcal{C}
  - \cot\theta\,(-\sin\theta\,Y_{,\theta})\,\mathcal{C}
  = \csc\theta\,Y_{,\phi\phi}\,\mathcal{C}
  + \cos\theta\,Y_{,\theta}\,\mathcal{C},
\end{equation}
where the second equality used $\Gamma^{\phi}_{\phi\theta}=\cot\theta$.
Adding,
\begin{equation}
  2\,\delta R_{\theta\phi}
  = -\bigl(\sin\theta\,Y_{,\theta\theta}
        - \csc\theta\,Y_{,\phi\phi}\bigr)\,\mathcal{C}.
\end{equation}
The angular factor
$\sin\theta\,Y_{,\theta\theta}-\csc\theta\,Y_{,\phi\phi}$ is in
general \emph{not} identically zero on $S^{2}$, so the full
projection $\delta R_{\theta\phi}=0$ requires
\begin{equation}
  \mathcal{C}(t,r) \;=\; 0
  \;\;\Longleftrightarrow\;\;
  \frac{\partial_{t}h_{0}}{f}
  - \partial_{r}\!\bigl(f h_{1}\bigr)
  + \frac{f h_{1}}{r}
  \;=\; 0,
\end{equation}
which, after multiplying through by $f$ and collecting,
is equivalent to
\begin{equation}
  \partial_{r}\!\bigl[f(r)\,h_{1}\bigr]
  - f(r)^{-1}\,\partial_{t} h_{0}
  \;=\; 0,
\label{eq:rw-C-derived}
\end{equation}
i.e.\ Eq.~\eqref{eq:rw-C}. (The $f h_{1}/r$ term in
$\mathcal{C}$ is reabsorbed into $\partial_{r}(f h_{1})$ via
the identity
$\partial_{r}(f h_{1}) - f h_{1}/r
  = (1/r)\,\partial_{r}\!\bigl[r f h_{1}\bigr] - 2 f h_{1}/r$,
or equivalently follows from rewriting
$\mathcal{C}$ as
$f\,\mathcal{C}=\partial_{t}h_{0}-f\,\partial_{r}(f h_{1})
  + f^{2} h_{1}/r$ and recognising the constraint as the
linearised form of the $r$-derivative of
Eq.~\eqref{eq:rw-Q-constraint} below.)

\paragraph{Linear-dependence checks.}
The components $\delta R_{t\phi}$ and $\delta R_{r\phi}$ produce
the same radial equations \eqref{eq:rw-A-derived} and
\eqref{eq:rw-B-derived} as $\delta R_{r\theta}$ and
$\delta R_{t\theta}$ respectively. Under the substitution
$(\csc\theta\,Y_{,\phi})\to(-\sin\theta\,Y_{,\theta})$ the
overall angular factor changes but the radial coefficients of
$h_{0},h_{1}$ are unchanged, since the axial ansatz
\eqref{eq:axial-components} pairs the two scalar bases
$(\csc\theta\,Y_{,\phi})$ and $(\sin\theta\,Y_{,\theta})$
symmetrically. The five axial components of $\delta R_{\mu\nu}$
therefore reduce to exactly three independent radial equations
in the two unknowns $(h_{0},h_{1})$: Eqs.~\eqref{eq:rw-A},
\eqref{eq:rw-B} and \eqref{eq:rw-C}.

\paragraph{Constraint preservation.}
The Bianchi identity $\nabla^{\mu}\delta G_{\mu\nu}=0$ on a
vacuum background, together with the fact that
\eqref{eq:rw-A-derived}--\eqref{eq:rw-B-derived} are evolution
equations for $(h_{0},h_{1})$ in $t$, implies that if
\eqref{eq:rw-C-derived} holds at $t=t_{0}$ then it holds for all
$t$, i.e.\ the constraint is propagated by the evolution
\eqref{eq:rw-A-derived}--\eqref{eq:rw-B-derived}. This is the
standard 3+1 statement that the Hamiltonian constraint of
\eqref{eq:rw-C-derived} is preserved by the time evolution of
the axial sector and need only be imposed on initial data.

This completes the explicit projection of the linearised vacuum
equations onto the axial harmonics. The reduction to a single
gauge-invariant scalar via the Regge--Wheeler function
\eqref{eq:rw-Q} and the master Eq.~\eqref{eq:rw-eqn} now
proceeds exactly as in the next paragraph.
